\documentclass[10pt,a4paper]{amsart}
\usepackage[margin=19mm]{geometry}
\usepackage{amsmath,amssymb,mathtools}
\usepackage[scr=rsfs,cal=euler]{mathalfa}
\usepackage{xfrac}
\usepackage[unicode,hidelinks,bookmarks=false]{hyperref}
\newcommand{\ie}{i.e.\ }
\newcommand{\cf}{cf.\ }
\newcommand{\resp}{respectively\ }
\newcommand{\Subsection}{\subsection}
\providecommand{\nobreakdash}{\nobreak\hbox{-}}
\setlength{\emergencystretch}{2em}
\multlinegap0pt
\usepackage{amssymb}
\usepackage{mathtools}
\usepackage{esint}
\usepackage[shortlabels]{enumitem}
\usepackage{tcolorbox}
\usepackage{float}
\usepackage{booktabs}
\usepackage{xcolor}
\usepackage{tensor}
\newtheorem{lemma}{Lemma}
\title{Asymptotic zero distribution of the polynomials $\widetilde{\Xi}_n$}
\author{Luc Rams\`es Talla Waffo}
\date{Source: arXiv:2602.20192v1 (CC BY 4.0)}
\begin{document}
\maketitle

\noindent\emph{Source and licence.} Asymptotic zero distribution of the polynomials $\widetilde{\Xi}_n$. Version arXiv:2602.20192v1 (CC BY 4.0), distributed by arXiv under the Creative Commons Attribution 4.0 International licence, http://creativecommons.org/licenses/by/4.0/, as recorded in the arXiv licence metadata for this exact version and shown on the abstract page https://arxiv.org/abs/2602.20192v1. Full licence notice: \texttt{licenses/arxiv-2602.20192-CC-BY-4.0.txt}.\par\smallskip

\noindent\textit{Editorial scope.} Counted unit: the article's lemma lemma (source0.tex lines 213--222). The article's complete original proof of it is delivered with it (source0.tex lines 224--366, one \texttt{proof} environment of 143 lines). No mathematics is written, repaired or re-derived: the statement and the proof are the article's own lines, in the article's own notation, and no retained line is substituted or re-flowed.  No further source-local context is needed: every cross-reference of every retained line resolves inside the two delivered spans.  Nothing was omitted from any retained span, so the package carries no omission record.  No retained line contains a citation, so the wrapper carries no bibliography and no bibliography entry.  No retained line contains a figure environment, an image-inclusion command or drawing code, so no image is delivered and the package declares no asset dependency.  Editorial formatting only, in addition: an AMS \texttt{amsart} preamble that restates the article's own declarations and macros used by the retained lines, namely the environments \texttt{lemma} on separate counters, together with the standard packages those lines need.  The retained environments print in this document as Lemma 1 in order of appearance, whereas the article prints the same items with the numbering of its own full document.  The complete native source of the article as distributed in the arXiv e-print for this exact version is bundled unchanged as \texttt{sources/195168/source0.tex}.
\begin{lemma}\label[lemma]{lemma:Sm-ratio}
Let
\[
S_m(x) := \sum_{k\ge 0} (2k+1)^m x^k,
\]
for $m\in\mathbb{N}$ and $0<x<1$. Then
\[
\lim_{m\to\infty} \frac{1}{m}\,\frac{S_{m+1}(x)}{S_m(x)} \;=\; \frac{2}{-\log x}.
\]
\end{lemma}

\begin{proof}
We first express $S_m(x)$ as a contour integral. For $n\in\mathbb{N}$ and $m\in\mathbb{N}$, Cauchy’s integral formula for derivatives gives
\[
n^m
= \left.\frac{d^m}{dz^m} e^{nz}\right|_{z=0}
= \frac{m!}{2\pi i} \oint_{|z|=r} \frac{e^{nz}}{z^{m+1}}\,dz,
\]
where $\{|z|=r\}$ is the positively oriented circle centered at $0$.

Hence
\[
(2k+1)^m
= \frac{m!}{2\pi i} \oint_{|z|=r} \frac{e^{(2k+1)z}}{z^{m+1}}\,dz.
\]
Substituting this into $S_m(x)$ and interchanging sum and integral (justified by absolute convergence for $|x|<1$ and $r$ sufficiently small) yields
\[
S_m(x)
= \frac{m!}{2\pi i} \oint_{|z|=r} \frac{1}{z^{m+1}}
   \sum_{k\ge 0} x^k e^{(2k+1)z}\,dz.
\]
The sum is geometric:
\[
\sum_{k\ge 0} x^k e^{(2k+1)z}
= e^z \sum_{k\ge 0} (x e^{2z})^k
= \frac{e^z}{1 - x e^{2z}},
\]
for $|x e^{2z}|<1$, which holds on a sufficiently small circle $|z|=r$. Thus
\[
S_m(x)
= \frac{m!}{2\pi i} \oint_{|z|=r}
  \frac{e^z}{z^{m+1} \bigl(1-x e^{2z}\bigr)}\,dz.
\]
Set
\[
g(z) := \frac{e^z}{1 - x e^{2z}},
\quad\text{so that}\quad
S_m(x) = \frac{m!}{2\pi i} \oint_{|z|=r} \frac{g(z)}{z^{m+1}}\,dz.
\]

We now analyze the singularities of $g$. The poles are determined by
\[
1 - x e^{2z} = 0
\quad\Longleftrightarrow\quad
e^{2z} = \frac{1}{x}.
\]
Since $0<x<1$, we have $1/x>1$, and all solutions are
\[
z_k = \frac{1}{2}\log\frac{1}{x} + \pi i k
    =: z_0 + \pi i k,\quad k\in\mathbb{Z},
\]
where
\[
z_0 := \frac{1}{2}\log\frac{1}{x} = -\frac{1}{2}\log x > 0
\]
is the unique pole of smallest modulus (real and positive). Each $z_k$ is a simple pole of $g$.

The residue of $g$ at $z_0$ is
\[
\operatorname{Res}(g,z_0)
= \frac{e^{z_0}}{(1 - x e^{2z})'(z_0)}
= \frac{e^{z_0}}{-2x e^{2z_0}}
= \frac{e^{z_0}}{-2x\cdot (1/x)}
= -\frac{e^{z_0}}{2}.
\]

Choose radii $0<r<|z_0|$ and $R$ such that
\(\displaystyle
|z_0| < R < \min_{k\ne 0} |z_k|.
\)
Enlarge the contour from the positively oriented circle $|z|=r$ to the positively oriented circle $|z|=R$. By the residue theorem,
\[
\oint_{|z|=r}\frac{g(z)}{z^{m+1}}\,dz
= \oint_{|z|=R}\frac{g(z)}{z^{m+1}}\,dz
  - 2\pi i\,\operatorname{Res}\Bigl(\frac{g(z)}{z^{m+1}}, z_0\Bigr),
\]
because in deforming the contour we enclose the additional pole at $z_0$ and no others (by the choice of $R$).

Since $g$ has a simple pole at $z_0$, we have
\(\displaystyle
\operatorname{Res}\Bigl(\frac{g(z)}{z^{m+1}}, z_0\Bigr)
= \frac{\operatorname{Res}(g,z_0)}{z_0^{m+1}}
= -\frac{e^{z_0}}{2 z_0^{m+1}}.
\)

Thus
\[
S_m(x)
= \frac{m!}{2\pi i} \oint_{|z|=r}\frac{g(z)}{z^{m+1}}\,dz
= \frac{m!}{2\pi i} \oint_{|z|=R}\frac{g(z)}{z^{m+1}}\,dz
  + m!\,\frac{e^{z_0}}{2 z_0^{m+1}}.
\]

On the positively oriented circle $|z|=R$, the function $g$ is analytic and bounded, say $|g(z)|\le M$ for all $|z|=R$. Hence
\(\displaystyle
\left|\frac{m!}{2\pi i} \oint_{|z|=R}\frac{g(z)}{z^{m+1}}\,dz\right|
\le \frac{m!}{2\pi} \cdot 2\pi R \cdot \frac{M}{R^{m+1}}
= m!\,M\,R^{-m}.
\)

Therefore
\[
S_m(x)
= m!\,\frac{e^{z_0}}{2 z_0^{m+1}}
  + O\bigl(m!\,R^{-m}\bigr),
\quad m\to\infty.
\]
The same argument with $m$ replaced by $m+1$ gives
\[
S_{m+1}(x)
= (m+1)!\,\frac{e^{z_0}}{2 z_0^{m+2}}
  + O\bigl((m+1)!\,R^{-(m+1)}\bigr).
\]

Thus
\(\displaystyle
\frac{S_{m+1}(x)}{S_m(x)}
= \frac{(m+1)!\,\dfrac{e^{z_0}}{2 z_0^{m+2}}
          \bigl(1 + O\bigl((R/z_0)^{-m}\bigr)\bigr)}
       {m!\,\dfrac{e^{z_0}}{2 z_0^{m+1}}
          \bigl(1 + O\bigl((R/z_0)^{-m}\bigr)\bigr)}
= \frac{m+1}{z_0}\bigl(1 + o(1)\bigr),
\quad m\to\infty,
\)
since $R>|z_0|$ implies $(R/z_0)^{-m}\to 0$.

Dividing by $m$ and letting $m\to\infty$ yields
\[
\lim_{m\to\infty} \frac{1}{m}\,\frac{S_{m+1}(x)}{S_m(x)}
= \lim_{m\to\infty} \frac{m+1}{m}\,\frac{1}{z_0}
= \frac{1}{z_0}.
\]
Because
\(\displaystyle
z_0 = \frac{1}{2}\log\frac{1}{x} = -\frac{1}{2}\log x,
\)
we obtain
\(\displaystyle
\frac{1}{z_0}
= \frac{2}{\log(1/x)}
= \frac{2}{-\log x}.
\)
This proves the claim.
\end{proof}

\end{document}
